%!TEX root = ../../thesis.tex
% Generated by scripts/evaluate.py -- do not edit; edit the emitter instead.
\begin{figure}[tbp]
  \centering
  \begin{subfigure}[b]{0.99\linewidth}
    \centering
    \begin{tikzpicture}
      \begin{axis}[thesis result plot, width=0.95\linewidth, height=4.2cm,
        xlabel={learner iteration}, ylabel={episode return},
        xmin=0, ymin=0, legend pos=south east, legend style={font=\scriptsize}]
      \addplot [draw=none, forget plot, name path=loA]
        table [col sep=comma, x=x, y=lo] {figures/data/return-cjson-ctrl-rand.csv};
      \addplot [draw=none, forget plot, name path=hiA]
        table [col sep=comma, x=x, y=hi] {figures/data/return-cjson-ctrl-rand.csv};
      \addplot [KITblue15, forget plot] fill between [of=loA and hiA];
      \addplot [KIT line plot A]
        table [col sep=comma, x=x, y=mid] {figures/data/return-cjson-ctrl-rand.csv};
      \addlegendentry{CTRL-RAND}
      \addplot [draw=none, forget plot, name path=loB]
        table [col sep=comma, x=x, y=lo] {figures/data/return-cjson-t25.csv};
      \addplot [draw=none, forget plot, name path=hiB]
        table [col sep=comma, x=x, y=hi] {figures/data/return-cjson-t25.csv};
      \addplot [KITred15, forget plot] fill between [of=loB and hiB];
      \addplot [KIT line plot B]
        table [col sep=comma, x=x, y=mid] {figures/data/return-cjson-t25.csv};
      \addlegendentry{T25}
      \end{axis}
    \end{tikzpicture}
    \caption{cJSON}
  \end{subfigure}
  \\[1.4ex]
  \begin{subfigure}[b]{0.99\linewidth}
    \centering
    \begin{tikzpicture}
      \begin{axis}[thesis result plot, width=0.95\linewidth, height=4.2cm,
        xlabel={learner iteration}, ylabel={episode return},
        xmin=0, ymin=0, legend pos=south east, legend style={font=\scriptsize}]
      \addplot [draw=none, forget plot, name path=loA]
        table [col sep=comma, x=x, y=lo] {figures/data/return-libxml2-libxml2-400.csv};
      \addplot [draw=none, forget plot, name path=hiA]
        table [col sep=comma, x=x, y=hi] {figures/data/return-libxml2-libxml2-400.csv};
      \addplot [KITblue15, forget plot] fill between [of=loA and hiA];
      \addplot [KIT line plot A, forget plot]
        table [col sep=comma, x=x, y=mid] {figures/data/return-libxml2-libxml2-400.csv};
      \end{axis}
    \end{tikzpicture}
    \caption{libxml2}
  \end{subfigure}
  \\[1.4ex]
  \begin{subfigure}[b]{0.99\linewidth}
    \centering
    \begin{tikzpicture}
      \begin{axis}[thesis result plot, width=0.95\linewidth, height=4.2cm,
        xlabel={learner iteration}, ylabel={episode return},
        xmin=0, ymin=0, legend pos=south east, legend style={font=\scriptsize}]
      \addplot [draw=none, forget plot, name path=loA]
        table [col sep=comma, x=x, y=lo] {figures/data/return-picohttpparser-picohttpparser-400.csv};
      \addplot [draw=none, forget plot, name path=hiA]
        table [col sep=comma, x=x, y=hi] {figures/data/return-picohttpparser-picohttpparser-400.csv};
      \addplot [KITblue15, forget plot] fill between [of=loA and hiA];
      \addplot [KIT line plot A, forget plot]
        table [col sep=comma, x=x, y=mid] {figures/data/return-picohttpparser-picohttpparser-400.csv};
      \end{axis}
    \end{tikzpicture}
    \caption{picohttpparser}
  \end{subfigure}
  \caption[Episode return against training iteration]{
    Per-trajectory return, median and interquartile range over the 128
    environments and four seeds.
    Return is the reward the agent optimises for and rises monotonically on all three targets.}
  \label{fig:return}
\end{figure}
